\documentclass[10pt,a4paper]{article}
\usepackage[T1]{fontenc}
\usepackage[utf8]{inputenc}
\usepackage{lmodern}
\usepackage{microtype}
\usepackage[margin=20mm,headheight=14pt,headsep=11pt,footskip=13mm]{geometry}
\usepackage{amsmath,amssymb,mathtools}
\usepackage{booktabs,tabularx,array}
\usepackage{xcolor}
\usepackage{enumitem}
\usepackage{listings}
\usepackage{hyperref}
\usepackage{fancyhdr}
\usepackage{titlesec}
\usepackage{needspace}

\definecolor{ink}{HTML}{172B47}
\definecolor{accent}{HTML}{16647B}
\definecolor{muted}{HTML}{576779}
\definecolor{pale}{HTML}{F2F6F8}
\definecolor{rulegray}{HTML}{D8E0E5}
\hypersetup{colorlinks=true,urlcolor=accent,linkcolor=accent,pdftitle={qaptiva-mbl User Guide},pdfauthor={qaptiva-mbl contributors},pdfsubject={AnalogQPU and exact diagonalization for disordered spin chains}}
\titleformat{\section}{\large\sffamily\bfseries\color{ink}}{\thesection}{0.65em}{}
\titleformat{\subsection}{\normalsize\sffamily\bfseries\color{accent}}{\thesubsection}{0.55em}{}
\titlespacing*{\section}{0pt}{10pt}{5pt}
\titlespacing*{\subsection}{0pt}{7pt}{3pt}
\setlength{\parindent}{0pt}
\setlength{\parskip}{4pt}
\renewcommand{\arraystretch}{1.18}
\setlist[itemize]{topsep=2pt,itemsep=2pt,parsep=0pt,leftmargin=*}
\setlist[enumerate]{topsep=2pt,itemsep=2pt,parsep=0pt,leftmargin=*}
\pagestyle{fancy}
\fancyhf{}
\fancyhead[L]{\scriptsize\sffamily\bfseries\color{accent} QAPTIVA-MBL}
\fancyhead[R]{\scriptsize\sffamily\color{muted} RESEARCH USER GUIDE}
\fancyfoot[L]{\scriptsize\sffamily\color{muted} Version 0.1.0 \;\textbar\; October 2026}
\fancyfoot[R]{\scriptsize\sffamily\color{muted} \thepage}
\renewcommand{\headrulewidth}{0.35pt}
\renewcommand{\footrulewidth}{0pt}
\lstdefinestyle{clean}{
  language=Python,basicstyle=\ttfamily\footnotesize,backgroundcolor=\color{pale},
  frame=single,framerule=0.3pt,rulecolor=\color{rulegray},
  numbers=none,xleftmargin=5pt,xrightmargin=5pt,framexleftmargin=5pt,
  framexrightmargin=5pt,breaklines=true,showstringspaces=false,
  columns=fullflexible,keepspaces=true,aboveskip=5pt,belowskip=7pt,
  keywordstyle=\color{accent}\bfseries,commentstyle=\color{muted},
  stringstyle=\color{ink}
}
\lstset{style=clean}
\newcommand{\ket}[1]{\lvert #1 \rangle}
\newcommand{\bra}[1]{\langle #1 \rvert}
\newcommand{\Udist}{\mathcal U}
\newcommand{\Bonds}{\mathcal B}
\newcommand{\Cq}{C^{\mathrm{QPU}}}
\newcommand{\Ce}{C^{\mathrm{ED}}}
\newcommand{\paulix}{\sigma^{x}}
\newcommand{\pauliz}{\sigma^{z}}

\begin{document}
\begin{center}
{\fontsize{22}{26}\selectfont\sffamily\bfseries\color{ink} qaptiva-mbl}\\[5pt]
{\large\sffamily\color{accent} Analog spin-chain dynamics and ED benchmarking}\\[8pt]
{\normalsize\color{muted} Minimal user guide \quad\textbar\quad v0.1.0 \quad\textbar\quad October 2026}
\end{center}
\vspace{4pt}{\color{rulegray}\hrule height 0.7pt}\vspace{10pt}

\section{Purpose and scope}
This package studies real-time dynamics in the \textbf{Imbrie} and \textbf{random-field XXZ} chains. A user selects physical parameters, boundary conditions and an initial product state; the package preserves the exact disorder realization, runs a Qaptiva \texttt{AnalogQPU} job, retrieves site-resolved observables and optionally compares the result against independent \textbf{exact diagonalization (ED)} using NumPy/SciPy.

\begin{tabularx}{\textwidth}{@{}>{\bfseries}p{38mm}X@{}}
\toprule
User setting & Available choices \\
\midrule
Model & \texttt{imbrie} or \texttt{xxz} \\
Geometry & Chain length $L$; \texttt{obc} or \texttt{pbc} \\
Initial state & N\'eel state or random $Z$-basis product state \\
Disorder & User-supplied seed and model couplings \\
QPU integration & Adaptive error control or a fixed number of steps \\
Outputs & Local autocorrelators, site average, ED comparison and error plots \\
\bottomrule
\end{tabularx}

The initial release supports \emph{one disorder realization per run}. There is no automatic disorder averaging, entanglement calculation or $U(1)$ symmetry-sector reduction.

\section{Hamiltonians and normalization}
Sites are labeled $j=0,1,\ldots,L-1$, and we set $\hbar=1$. Let $X_j=\sigma_j^x$, $Y_j=\sigma_j^y$, and $Z_j=\sigma_j^z$ be Pauli operators. Neighbor pairs are
\begin{equation}
\Bonds_{\mathrm{OBC}}=\{(j,j+1)\mid 0\le j<L-1\},\qquad
\Bonds_{\mathrm{PBC}}=\Bonds_{\mathrm{OBC}}\cup\{(L-1,0)\}.
\end{equation}
OBC requires $L\ge2$; PBC requires $L\ge3$ in this implementation.

\subsection{Imbrie model: Pauli convention}
\begin{equation}
\boxed{H_{\mathrm I}=\sum_{(i,j)\in\Bonds}\Delta_{ij}Z_iZ_j
+\sum_{j=0}^{L-1}\big(h_jZ_j+\Gamma X_j\big).}
\end{equation}
The bond couplings and longitudinal fields are independent draws:
\begin{equation}
\Delta_{ij}\sim\Udist(\Delta_{\min},\Delta_{\max}),\qquad
h_j\sim\Udist(-h_0,h_0).
\end{equation}
The Hamiltonian uses Pauli matrices \emph{directly}, without extra factors of $1/2$. Configuration keys: \texttt{gamma}, \texttt{delta\_min}, \texttt{delta\_max}, \texttt{h0}.

\subsection{Random-field XXZ: spin-\texorpdfstring{$\tfrac12$}{1/2} convention}
With $S_j^\alpha=\sigma_j^\alpha/2$ and $h_j\sim\Udist(-h_0,h_0)$,
\begin{align}
H_{\mathrm{XXZ}}={}&\sum_{(i,j)\in\Bonds}\!\big[J_{xy}(S_i^x S_j^x+S_i^y S_j^y)+J_z S_i^zS_j^z\big]
+\sum_jh_jS_j^z \\[-2pt]
={}&\sum_{(i,j)\in\Bonds}\!\left[\frac{J_{xy}}{4}(X_iX_j+Y_iY_j)
+\frac{J_z}{4}Z_iZ_j\right]+\frac12\sum_jh_jZ_j.\notag
\end{align}
The \emph{second line} specifies the exact Pauli coefficients used for both solvers. Configuration keys: \texttt{jxy}, \texttt{jz}, \texttt{h0}.

\clearpage
\section{Initial states and measured observables}
The package prepares a product eigenstate of all $Z_j$,
\begin{equation}
\ket{\psi_0}=\ket{s_0,s_1,\ldots,s_{L-1}},\qquad
Z_j\ket{\psi_0}=s_j\ket{\psi_0},\qquad s_j\in\{+1,-1\}.
\end{equation}
\begin{itemize}
\item \textbf{N\'eel:} \texttt{initial\_state="neel"}, $s_j=(-1)^j$.
\item \textbf{Random:} \texttt{initial\_state="random"}, independent $s_j=\pm1$.
\end{itemize}
In the Qaptiva input bitstring, $s_j=+1$ maps to \texttt{0} and $s_j=-1$ to \texttt{1}. \textbf{Site 0 is the leftmost bit.} Independent deterministic pseudorandom streams separate Hamiltonian disorder from random-state selection.

For a time-independent Hamiltonian, $\ket{\psi(t)}=e^{-iHt}\ket{\psi_0}$. Define $Z_j(t)=e^{iHt}Z_je^{-iHt}$. The product-state site-resolved autocorrelator is
\begin{equation}
\boxed{C_j(t)=\bra{\psi_0}Z_j(t)Z_j\ket{\psi_0}
=s_j\bra{\psi(t)}Z_j\ket{\psi(t)},}
\qquad
\boxed{C(t)=\frac1L\sum_{j=0}^{L-1}C_j(t).}
\end{equation}
The QPU job measures unweighted local magnetizations $\langle Z_j(t)\rangle$. After retrieval, the package multiplies each site by the \emph{stored} initial spin $s_j$. For ideal unitary dynamics $C_j(0)=1$.

\textbf{Interpretation:} for one product state, $C_j(t)$ is \emph{not} the infinite-temperature trace correlation $2^{-L}\operatorname{Tr}[Z_j(t)Z_j]$; the latter requires a trace or suitable state averaging.

\section{Solver settings: adaptive or fixed-step}
The \texttt{RunConfig} constructor enforces \textbf{mutually exclusive} integration settings:

\begin{tabularx}{\textwidth}{@{}p{30mm}X@{}}
\toprule
\textbf{Adaptive} & \texttt{error\_control=True}, \texttt{nsteps=None}; provide positive \texttt{abs\_tol} and \texttt{rel\_tol}. Qaptiva selects internal steps through its adaptive Dormand--Prince Runge--Kutta (DP5) method. \\
\midrule
\textbf{Fixed steps} & \texttt{error\_control=False}; supply a positive integer \texttt{nsteps}. Adaptive tolerances are not forwarded to the remote solver. \\
\bottomrule
\end{tabularx}

\begin{lstlisting}
# Adaptive (default): do not specify nsteps
error_control=True, nsteps=None, abs_tol=1e-8, rel_tol=1e-8

# Fixed-step alternative
error_control=False, nsteps=1000
\end{lstlisting}

The \emph{internal integration steps} are not the same as the \emph{returned observation times}. The package always takes timestamps from the QPU result. The \texttt{use\_gpu} setting controls the backend's GPU option where available.

\section{Notebook as the main interface}
Run \texttt{notebooks/MBL\_Simulation.ipynb}. Its cells perform configuration, sample preparation, optional submission, optional retrieval, ED evolution, plotting and comparison. The remote cells start \emph{disabled}; enabling them requires a configured Qaptiva installation.

\section{Installation}
Requires \textbf{Python 3.10+}, NumPy, SciPy and Matplotlib. Run these commands from the repository root:
\begin{lstlisting}[language=bash]
python -m pip install -e ".[notebook,test]"
python -m pytest -q
jupyter notebook notebooks/MBL_Simulation.ipynb
\end{lstlisting}
The Qaptiva/QLMaaS modules \texttt{qat} and \texttt{qlmaas}, together with authenticated server access, are required only for remote submission/retrieval; they are not bundled in this repository. ED-only workflows run locally.

\clearpage
\section{Prepare, submit and retrieve}
Use the notebook, or call the public API:
\begin{lstlisting}
from qaptiva_mbl import RunConfig, prepare_run
from qaptiva_mbl import submit_run, retrieve_run

cfg = RunConfig(
    model="imbrie", L=8, boundary="pbc", seed=12345,
    initial_state="neel", gamma=1.0,
    delta_min=0.8, delta_max=1.2, h0=3.5,
    tmax=100.0, error_control=True, nsteps=None,
    abs_tol=1e-8, rel_tol=1e-8, use_gpu=False,
)
run_dir = prepare_run(cfg, root="runs")
job_id = submit_run(run_dir)   # requires Qaptiva access
print(job_id)

# Retrieve now or in a later session from the same run_dir.
times, C_qpu = retrieve_run(run_dir, host="YOUR-QLMAAS-HOST")
\end{lstlisting}

For the XXZ chain, change \texttt{model="xxz"} and specify \texttt{jxy} and \texttt{jz}. The package persists the random fields, Imbrie bonds (if present), initial spins, parameters and remote job ID. It never needs to redraw a disorder sample at retrieval time.

\section{Independent manual ED}
ED constructs the same spin Hamiltonian as a sparse matrix in the \emph{full} $2^L$-dimensional basis, without using Qaptiva's matrix constructors. The time evolution uses
\begin{equation}
\ket{\psi(t)}=e^{-iHt}\ket{\psi_0}.
\end{equation}
The \texttt{dense} method diagonalizes $H$ once, while \texttt{krylov} applies sparse matrix exponentials. With \texttt{method="auto"}, dense ED is used for $L\le9$ and Krylov for larger $L$. The default \texttt{max\_L=12} safety limit can be raised explicitly, at increased computational cost.

\begin{lstlisting}
from qaptiva_mbl import run_ed, compare_run

# Reuse the exact QPU timestamps, saved disorder and initial state.
times, C_ed = run_ed(run_dir, max_L=12)
metrics = compare_run(run_dir)
print(metrics["global_max_absolute_error"])
\end{lstlisting}

For \textbf{ED without Qaptiva}, prepare a run and supply a time array:
\begin{lstlisting}
import numpy as np
from qaptiva_mbl import RunConfig, prepare_run, run_ed
run_dir = prepare_run(RunConfig(model="xxz", L=8))
times, C_ed = run_ed(run_dir, times=np.linspace(0, 100, 201))
\end{lstlisting}

\clearpage
\section{Quantifying the ED--AnalogQPU difference}
For the \emph{same} Hamiltonian realization, initial state and returned timestamps $t_k$, define
\begin{equation}
\delta C_j(t_k)=C_j^{\mathrm{QPU}}(t_k)-C_j^{\mathrm{ED}}(t_k),\qquad
\epsilon_j(t_k)=|\delta C_j(t_k)|.
\end{equation}
The analysis also computes the site-resolved average difference and three nonnegative aggregate errors:
\begin{align}
\delta C(t_k)&=\frac1L\sum_{j=0}^{L-1}\delta C_j(t_k),\\
\mathrm{MAE}(t_k)&=\frac1L\sum_{j=0}^{L-1}|\delta C_j(t_k)|,\\
\mathrm{RMSE}(t_k)&=\sqrt{\frac1L\sum_{j=0}^{L-1}[\delta C_j(t_k)]^2},\\
\epsilon_{\max}(t_k)&=\max_{0\le j<L}|\delta C_j(t_k)|.
\end{align}
A global maximum, $\max_{j,k}|\delta C_j(t_k)|$, global RMSE and per-site maximum errors are also saved. Note that positive and negative deviations can cancel in $\delta C(t)$; it is not interchangeable with MAE. These metrics quantify \emph{solver agreement}, not extrapolation to the thermodynamic limit.

\section{Plots and results directory}
After retrieving QPU data, figures can be regenerated without any network call:
\begin{lstlisting}
from qaptiva_mbl import plot_run
plot_run(run_dir, kind="sites", sites=[0,1,2,3],
         save="sites.pdf")
plot_run(run_dir, kind="average", save="average.pdf")
plot_run(run_dir, kind="error", save="errors.pdf")
plot_run(run_dir, kind="site_error", sites=[0,1],
         save="site_errors.pdf")
\end{lstlisting}
Optional arguments include \texttt{log\_time=True} and \texttt{overlay\_ed=True}. The function returns a Matplotlib figure for customization. In ED-only mode, the notebook plots the saved ED arrays directly.

\begin{lstlisting}[language={}]
runs/imbrie_L8_<unique-id>/
    run.json          # run configuration, bitstring, job ID, status
    sample.npz        # exact fields, bond couplings, initial spins
    qpu_results.npz   # time, local magnetization and correlators
    ed_results.npz    # time and ED correlators
    comparison.npz    # signed/absolute differences and summaries
\end{lstlisting}
Only files produced by completed stages exist. The result arrays are indexed as $[j,k]$ for site $j$ and observation time $t_k$.

\section{Reliability notes and status}
\begin{itemize}
\item \textbf{Time alignment:} comparison rejects mismatched observation grids. ED must use the stored QPU times.
\item \textbf{Observable order:} site 0 is read from \texttt{value\_data}; remaining sites are read from \texttt{values\_data} in order.
\item \textbf{Reproducibility:} full-precision random arrays and the initial bitstring are saved before submission; job metadata is retained in \texttt{run.json}.
\item \textbf{Security:} never store credentials in source, notebooks, run metadata or Git.
\item \textbf{Validation status:} local tests and mock retrieval pass; check live Qaptiva submission and retrieval against the intended server.
\end{itemize}

\noindent\textbf{References:} \href{https://qlm.bull.com/qlm/notebooks/tutorials/analog_simulations/overview.html}{Qaptiva analog simulation overview} \;\textbar\;
\href{https://qlm.bull.com/qlm/notebooks/tutorials/analog_simulations/aqpu_getting_started.html}{Getting started} \;\textbar\;
\href{https://qlm.bull.com/qlm/notebooks/tutorials/schedule_arithmetic.html}{Schedules}.\\
\textbf{Project:} \url{https://github.com/abdpdn/qaptiva-mbl}

\end{document}
